% a01_H0rel part 1: setting on the cube, the weight class, the photon sector
\subsection{The cube in block units, and the weighted block average}\label{app:a01:weights}

\paragraph{Block units.} Fix the spacing index $k$ and a step $j$, $0\le j\le k$, and write everything in \emph{unit-$j$ units}: the $j$-blocks (cubes of side $5^{j}\eps_k=\eta$ in physical units) have side one, and the fine lattice has spacing $\vth\defeq5^{-j}$, with $1/\vth\in5\N$ for $j\ge1$. The number of blocks per side is $N\defeq2MJ/\eta\in\N$ (the side $2MJ$ of $\CJ$ is a multiple of every block side, \cref{ssec:k01:setting}). After a translation the picture is: fine sites $\Lambda\defeq\vth\Z^3$, the \emph{killed set} $\Kill\defeq\{x\in\Lambda\cap[0,N]^3:\ x_\mu\in\{0,N\}\text{ for some }\mu\}$ (the two Dirichlet planes of $\CK$ in each direction, $2MJ$ apart), the \emph{live set} $\Xl\defeq\{x\in\Lambda:\ 0<x_\mu<N\ \forall\mu\}$, which is $\alive^K_J$ of \cref{def:s02:Kcube} ($N/\vth-1$ sites per side), the blocks $B(y)\defeq(y+[0,1)^3)\cap\Lambda$, $y\in\Z^3$, $0\le y_\mu\le N-1$, and the \emph{face layer} $\Lone$ of blocks with some $y_\mu\in\{0,N-1\}$. The bottom Dirichlet plane $x_\mu=0$ is the first site plane of the blocks with $y_\mu=0$ (inside those blocks); the top plane $x_\mu=N$ is the first site plane of the blocks $y_\mu=N$, which lie outside the cube. We write $\DD$ for the union of the Dirichlet planes, $t(x)\defeq\min_\mu\min(x_\mu,N-x_\mu)$ for the face distance of a live site ($t=\vth$ on the first live layer $F_1$, the live sites with a neighbour in $\Kill$), and a \emph{face bond} is a bond from a live site to a killed site. By \cref{def:s02:acube}, a face bond carries only the diagonal term $\frac12\vth^{d-2}|\phi(x)|^2$ at its live end: it has no link variable, and $D_A\phi(b)=-\vth^{-1}\phi(x)$ on it whatever $A$. The photon is a site field on $\XJ$ with Neumann conditions on the bond planes $x_\mu=-\vth/2$, $N-\vth/2$ (\cref{def:s02:Kcube}).

\paragraph{The weight.} King's Dirichlet weight \cite[(3.8)--(3.9)]{K2}, in the form in which the averaging blocks have their top layers on the Dirichlet planes (\cref{ssec:k01:setting}), is the function
\[
w(x)\ \defeq\ \Wt\prod_{\mu=1}^{3}\mathbf 1\big[x_\mu/\vth\not\equiv0\ (\mathrm{mod}\ 5)\big],\qquad \Wt\defeq(5/4)^3=125/64,
\]
on the fine sites of a block: zero on the three site planes of the block that lie in the residue class of the Dirichlet planes (the \emph{zero-weight layers}), the constant $\Wt$ elsewhere. Since $1/\vth\equiv0\pmod5$, the block corners and all of $\Kill$ lie in the zero-weight class, and $\sum_{x\in B(y)}\vth^3w(x)=1$ for every block of the cube. The weight has period $5\vth$, so it acts at the finest step only: the coarse averages above scale one are uniform, $Q_{j+1}=Q\,Q_j$ with $Q$ the uniform one-step average with $QQ^*=1$, and the constants $a_j$ of \cite[(2.13)]{K1} are those of the uniform flow. The block average and its adjoint are
\[
Q(A)\phi(y)\ \defeq\ \sum_{x\in B(y)\cap\Xl}\vth^3\,w(x)\,U(A(\Gamma_{y,x}))\,\phi(x),\qquad P(A)\defeq Q(A)^*Q(A),
\]
with King's hierarchical contours $\Gamma_{y,x}$ \cite[(2.10)--(2.12)]{K1}, chosen coordinate-monotone between the block centre and $x$; such a path between two live sites of one block meets no killed site, since the killed sites of a block sit at the extreme value of a coordinate. The adjoint is taken for the measure $\vth^3$ on fine sites and counting measure on blocks.

\begin{definition}[weight class]\label{def:a01:wclass}
A block weight is a function $w\ge0$ on the fine sites of a unit block with (W1) $\sum_{x\in B}\vth^3w(x)=1$; (W2) $w\le\Wt$; (W3) $w=0$ on every killed site of the block; (W4) the same $w$ on every block up to translation. The uniform weight ($w\equiv1$, $\Wt=1$) and King's weight above ($\Wt=125/64$) are in the class.
\end{definition}

\begin{lemma}[the weighted average]\label{lem:a01:WQ}
For every weight in the class and every gauge field $A$, with $W_2\defeq\sum_{x\in B}\vth^3w(x)^2\in[1,\Wt]$ (for King's weight $W_2=\Wt$):
\begin{enumerate}[label=(Q\arabic*),nosep]
\item $Q(0)\mathbf 1=\mathbf 1$ on every block, and $Q(A)$ maps the transported constant to the constant;
\item $Q(A)Q(A)^*=W_2\,\mathbf 1$;
\item $P=W_2\Pi$ with $\Pi$ an orthogonal projector; $0\le P\le W_2$ and $\norm{1-P}=\max(1,|1-W_2|)\le1$ since $W_2\le\Wt<2$;
\item $|(Qu)(y)|\le W_2^{1/2}\norm{\mathbf 1_{B(y)}u}_2$ and $\norm{Q^*e_y}_2=W_2^{1/2}$;
\item $|(Qu)(y)|\le\sup_{B(y)}|u|$;
\item $|Q^*b|_\infty\le\Wt|b|_\infty$, $\norm{P}_{\infty\to\infty}\le\Wt$, and the Schur norm $\sup_x\sum_{x'}\vth^3|P(x,x')|\le\Wt$;
\item $\sum_x\vth^3|(Pu)(x)|=\sum_y|(Qu)(y)|$;
\item $P(x,x')=0$ unless $x,x'$ lie in one block, and $P$ vanishes on killed sites: $P$ never couples a live site to a killed one, and every union of blocks is invariant;
\item under $\phi\mapsto e^{eq\lambda}\phi$, $A\mapsto A+d\lambda$, the operators $Q(A)$, $P(A)$ transform covariantly.
\end{enumerate}
\end{lemma}
\begin{proof}
(Q1), (Q2): (W1) and the orthogonality of the transports $U$. (Q3): $P^2=Q^*(QQ^*)Q=W_2P$, and $1-P$ has the eigenvalues $1$ on $\Pi^\perp$ and $1-W_2$ on $\Pi$. (Q4): Cauchy--Schwarz with the weight $w$ and (W1); $\norm{Q^*e_y}^2=\sum\vth^3w^2$. (Q5): $|Qu|\le\sum\vth^3w|u|\le\sup|u|$ by (W1). (Q6): $(Q^*b)(x)=w(x)U^*b(y)$ and $(Pu)(x)=w(x)\sum_{x'}\vth^3w(x')U\,u(x')$. (Q7): $|(Pu)(x)|=w(x)|(Qu)(y)|$ and (W1). (Q8): the definitions and (W3). (Q9): the contour transports transform with the endpoint gauge functions, and $w$ is a scalar.
\end{proof}

Where a step of Ba{\l}aban's or King's torus analysis uses $QQ^*=1$, $P^2=P$, or that the block average of a constant is that constant, the cube version below uses the corresponding item of \cref{lem:a01:WQ} and acquires a factor depending on $\Wt$ only. Three identities change in substance and are proved below: the background of a constant datum (\cref{lem:a01:RC}), the block-mismatch identity (\cref{lem:a01:EF}), and the minimiser equation with $QQ^*=W_2$ (used in \cref{lem:a01:EB}).

\paragraph{The operators.} In unit-$j$ units let $V_j\defeq m^2\eta^2\le c_V$ be the mass term, $A$ a real gauge field on the live bonds of $\Xl$ with $U(A_b)=\exp(\vth\,e_jqA_b)$, $e_j\defeq e\eta^{1/2}$, and
\[
\Mop(\Omega,A)\ \defeq\ -\Delta^{N,D}_{A,\Omega}+V_j+a_jP(A)\qquad\text{on }\ell^2(\Omega\cap\Xl),
\]
for a \emph{sub-domain} $\Omega$ (a union of $j$-blocks of the cube; $\Omega=\Xl$ is the whole cube): the Laplacian keeps every bond with both ends in $\Omega\cap\Xl$, adds $\vth^{-2}$ to the diagonal for each face bond (Dirichlet, no link), and on bonds from $\Omega\cap\Xl$ to $\Xl\setminus\Omega$ either drops them (Neumann cut) or adds the diagonal (Dirichlet cut). The \emph{Neumann boundary} $\dN\Omega$ is the set of block midplanes between a block of $\Omega$ and a block of the cube not in $\Omega$; the \emph{depth} of a point is $\dep(x)\defeq\dist_\infty(x,\dN\Omega)$ ($+\infty$ on the whole cube). The Dirichlet faces of the cube are part of the operator, never part of $\dN\Omega$: every block touching a Dirichlet plane contains that plane's killed sites (bottom) or ends one site before it (top), so no sub-domain has a Neumann face on a Dirichlet plane. The \emph{leg} of block data $b$ is $u\defeq a_j\Mop(\Omega,A)^{-1}Q(A)^*b$; the \emph{background} of the block field $\phi_j$ is $\phi^{(j)}\defeq a_j\Mop^{-1}Q^*\phi_j$ \cite[(2.14)--(2.16)]{K1}; the block form is $\Delta^{(j)}\defeq a_j-a_j^2Q\Mop^{-1}Q^*$ and the fluctuation covariance $C^{(j)}\defeq(\Delta^{(j)}+aL^{-2}Q_c^*Q_c)^{-1}$ with $Q_c$ the uniform coarse average. A field $A$ is \emph{regular} on $\Omega$ at scale $j$ if it obeys \cite[(3.5)]{K1}: $|\partial_\mu A_\nu|\le K_j$ with $\theta_j\defeq e_jK_j\le\theta_0$ small; King's small-field class gives $\theta_j\le Ce_jp_j$, small for $\eta\le\eta_1$. Gauge covariance on the cube: for any real $\lambda$ on the sites, $\Mop(\Omega,A+d\lambda)=U_\lambda\Mop(\Omega,A)U_\lambda^*$ with $U_\lambda$ the multiplication by $e^{eq\lambda}$, $Q(A+d\lambda)=U^c_\lambda Q(A)U_\lambda^*$, and on a face bond $|D_{A+d\lambda}(U_\lambda u)|=|D_Au|=\vth^{-1}|u(x)|$; so every bound at a constant field $B=d(B\cdot x)$ on a box equals the bound at $0$ for a datum of the same modulus (the contours lie in the block, the Dirichlet condition is multiplication-invariant, the weight is a scalar).

\paragraph{The small-field set and the face test.} The small-field indicator $\chi_j$ of \cref{ssec:k01:setting} is \cite[(3.2)]{K1} on the blocks of the cube, with the gradient clause $|D_{A_j}\phi_j(\beta)|\le p_{j-1}$ on every block bond $\beta$ of the cube (all of which are cube-internal), times King's face test \cite[(5.41)]{K2}
\[
\lambda_j\ \defeq\ \prod_{y'\in\Lone}\chi\big(|\phi_j(y')|\le\sigma_j\big),\qquad \sigma_j\defeq\Big(\frac{2}{\theta_2\gamma_D}\Big)^{1/2}p_j\ \le\ C_8\,p_j,\quad \gamma_D\defeq\frac{a_j}{1+a_j\Wt^2},
\]
with $\theta_2\in(0,1)$ a fixed share and $\gamma_D$ the Dirichlet-face positivity constant of \cref{lem:a01:FPC}; $p_j\defeq p(\eta)=b_0(1+\ln\eta^{-1})^{p}$ and $M_j\defeq p_{j-1}(\lambda\ell_{j-1})^{-1/4}$ is the block-sup threshold of \cite[(3.2)]{K1} ($\ell_j\defeq5^j\eps_k$ the physical block side, $\lambda(\ell_j)=\lambda\ell_j$ the running coupling in $d=3$). The constant $C_8$ depends on $(a,\Wt,\theta_2,L)$ only.

\subsection{The photon sector: even restriction of the doubled torus}\label{app:a01:photon}

Let $2X$ be the torus of period $2N$ in each direction (physical period $4MJ$), $\XJ$ embedded as $\{0,\dots,N-\vth\}^3$ after the translation above, and $\mathcal W\cong(\Z_2)^3$ the group generated by the reflections $r_\mu$ in the photon planes $x_\mu=-\vth/2$ (on $2X$ the reflection in $x_\mu=N-\vth/2$ is the same element). Since $4MJ\equiv0$ modulo every block side, the blocks of $\XJ$ are blocks of the periodic tiling of $2X$, and $\mathcal W$ maps blocks to blocks. For a photon component $A_\mu$ (a site field) let $\Ebar$ be the even extension to $2X$, and $\Ebar_c$ the same on the coarse lattices.

\begin{lemma}[photon objects are even restrictions]\label{lem:a01:PE}
For every photon object $O$ of Ba{\l}aban's and King's construction built from the Neumann form $-\Delta^N_{\XJ}+\lmu^2\eta^2$ and the uniform block averages on $\XJ$, and its counterpart $O_{2X}$ on the periodic torus $2X$:
\begin{enumerate}[label=(\alph*),nosep]
\item the propagators $G_j$, the block forms $\Delta^{(j)}$, the fluctuation covariances $C^{(j)}$, the background map $A_j\mapsto A^{(j)}=a_jG_jQ_j^*A_j$ and the slices satisfy $O_{2X}\Ebar=\Ebar\,O$;
\item quadratic forms: $\ip{\Ebar A}{\Delta^{(j)}_{2X}\Ebar A}=8\ip{A}{\Delta^{(j)}A}$, and likewise for every positive form above;
\item for every coarse $\Lambda\subset X^{(j)}$ (resp.\ every fine union of blocks $\Omega\subset\XJ$) the Dirichlet-conditioned compressed operators on $\Lambda$ and the Neumann operators on $\Omega$ intertwine with $\Ebar$ and the $\mathcal W$-invariant sets $\mathcal W\Lambda$, $\mathcal W\Omega$;
\item if Ba{\l}aban's cut-off kernel $\zeta$ of \cite[(2.44)]{Ba2} is applied image-wise on $2X$ (the cube kernel being $K_X(x,y)\defeq\sum_{g\in\mathcal W}\zeta(x,gy)K_{2X}(x,gy)$), the cut-off objects of the cube are even restrictions of those of $2X$ and the support and size properties of $\zeta$ hold for them with the torus constants;
\item every positivity bound of the photon sector (\cite[(3.29), Prop.~3.1]{Ba2} photon part, the last-step mass bound $G'_k\ge\gamma_1\lmu^2\ell_k^2$ of \cite[(3.11)]{Ba2}, the photon exponents of \cite[(3.26)]{Ba2}) holds on $\XJ$ with the constant of $2X$.
\end{enumerate}
\end{lemma}
\begin{proof}
(a) The Neumann Laplacian on $\XJ$ is the periodic Laplacian of $2X$ on $\mathcal W$-even fields with reflection planes between sites: every plane-crossing bond of $2X$ joins a site to its mirror image, where an even field has zero difference, and every bond of $\XJ$ appears $8$ times. Block averages commute with $\Ebar$ (blocks go to blocks) and $Q^*_{2X}\Ebar_c=\Ebar Q^*_X$ (the adjoint of a $\mathcal W$-covariant uniform average is the uniform block extension). Hence $(M_{2X}+aQ^*Q)\Ebar=\Ebar(M_X+aQ^*Q)$; the torus operator commutes with $\mathcal W$, preserves the even subspace and is invertible on it, so $G_{2X}\Ebar=\Ebar G_X$, and $\Delta^{(j)}$, $C^{(j)}$, the background map and the slices follow from the algebra of \cite[(2.13)--(2.17)]{K1}, which is the same on both sides. (b) $\Ebar^*\Ebar=8$ on $\XJ$. (c) Compression to a $\mathcal W$-invariant set commutes with $\mathcal W$, and $P_{\mathcal W\Lambda}\Ebar=\Ebar P_\Lambda$; then (a) on the invariant subspace. For Neumann conditions on $\Omega$, the bonds of $\mathcal W\Omega$ are the images of those of $\Omega$ plus plane-crossing bonds, on which even fields have zero difference. (d) $\zeta(x,gy)$ depends on $|x-gy|$ only; the image-wise kernel commutes with $\mathcal W$ by construction and restricts to the cube kernel by definition; the support and size statements are statements about $\zeta(x,\cdot)$ on $2X$, and $|x-gy|\ge|x-y|$ for $x,y\in\XJ$ keeps every support statement. The cut is a device inside the definition of Ba{\l}aban's background, so the choice is free. (e) By (b): the cube form of an even field is $1/8$ of the torus form, and the right-hand sides (gradient sums, mass sums, block sums) scale by the same factor since plane-crossing gradients of even fields vanish.
\end{proof}

\begin{lemma}[regularity is preserved]\label{lem:a01:REG}
(a) A block photon configuration $A_j$ on $\XJ$ satisfies the photon clauses of \cite[(3.2)]{K1} if and only if $\Ebar_cA_j$ does on $2X$. (b) If $A$ on $\XJ$ is regular in the sense of \cite[(3.5)]{K1}, or obeys a H\"older bound with constant $C$, so does $\Ebar A$ on $2X$ with the same constant.
\end{lemma}
\begin{proof}
(a) Plane-crossing coarse gradients of $\Ebar_cA_j$ vanish; the other clauses are copies. (b) $\Ebar A=A\circ F$ with $F$ the coordinatewise folding map $2X\to\XJ$, which is $1$-Lipschitz; lattice differences across a plane vanish, the others are copies.
\end{proof}

\begin{lemma}[Neumann photon lines are image sums]\label{lem:a01:NL}
Every photon line of the cube (the fluctuation covariances and slices of \cref{lem:a01:PE}(a) and the full covariance $(-\Delta^N+\lmu^2\eta^2)^{-1}$) satisfies $K_X(x,y)=\sum_{g\in\mathcal W}K_{2X}(x,gy)$ for $x,y\in\XJ$. Hence every line bound of the torus analysis (decay, pointwise and derivative clauses, H\"older clauses in either variable, moment forms) holds for $K_X$ with the constant multiplied by $8$ and the same rates.
\end{lemma}
\begin{proof}
Kernel form of \cref{lem:a01:PE}(a): $(K_Xf)(x)=(K_{2X}\Ebar f)(x)=\sum_y\sum_gK_{2X}(x,gy)f(y)$. Each image term obeys the torus bound at the pair $(x,gy)$, and $|x-gy|\ge|x-y|$; derivatives in $x$ act on $K_{2X}(\cdot,gy)$, derivatives in $y$ through the lattice isometry $y\mapsto gy$.
\end{proof}

Consequently every photon statement that the torus analysis takes from \cite{Ba1,Ba2,Ba3} holds on $\XJ$: apply it on the torus $2X$ (whose period $4MJ$ is a multiple of every large-block side, and whose sides exceed the diameter of every box used) to the even data and restrict. For $x\in\XJ$, $\dist_{2X}(x,\partial(\mathcal W\Omega))\ge\dist(x,\partial\Omega\setminus\partial\XJ)$, so the restricted-form clauses transfer, and at the faces of $\XJ$ the photon has no boundary at all in the doubled picture. The photon tadpole and exchange graphs of the counterterm that contain a Neumann line $C_N-C_{\mathrm{free}}=\sum_{g\ne1}C_{\mathrm{free}}(\cdot,g\,\cdot)$ are bounded by the same image argument: with $\varrho(x)$ the distance of $x$ to the nearest photon plane, $|x-gx|\ge2\varrho(x)$ for $g\ne1$ and $|C_{\mathrm{free}}(x,z)|\le C\min(\vth^{-1},|x-z|^{-1})e^{-\lmu|x-z|}$, so the photon-image part of the mass counterterm is $O\big((e^2+e^4)(\min(\vth^{-1},\varrho^{-1})+1)\big)$ at order $\le4$, the same size as the scalar boundary remainders near the face, and it is absorbed in the zone-B bounds of \cref{prop:a01:TD} (it is the photon-image part of the boundary remainder families there).
